\documentclass[12pt]{article}
\usepackage{amsmath, amssymb, amsthm}
\usepackage{array}
\usepackage{booktabs}
\usepackage[margin=1in]{geometry}
\setlength{\headheight}{14.5pt}
\addtolength{\topmargin}{-2.5pt}
\usepackage{xcolor}
\usepackage{tcolorbox}
\tcbuselibrary{breakable}
\usepackage{enumitem}
\usepackage{fancyhdr}
\usepackage{graphicx}

% Define solution environment
\newtcolorbox{solution}{
    colback=blue!5!white,
    colframe=blue!75!black,
    title=Solution,
    fonttitle=\bfseries,
    breakable
}

\newtcolorbox{explanation}{
    colback=green!5!white,
    colframe=green!75!black,
    title=Explanation,
    fonttitle=\bfseries,
    breakable
}

\title{CHAPTER-WISE PROBLEMS\\
\large Regression Analysis}
\author{Statistical Foundation of Data Science\\
CSU1658}
\date{December 2025}

\pagestyle{fancy}
\fancyhf{}
\rhead{Regression Problems}
\lhead{CSU1658}
\cfoot{\thepage}

\begin{document}

\maketitle

\tableofcontents
\newpage

% ============================================================
% LINEAR REGRESSION PROBLEMS
% ============================================================

\section{Linear Regression}

\subsection*{Problem 1: Simple Linear Regression with Hypothesis Testing}

A researcher studies the relationship between study hours (\(x\)) and exam scores (\(y\)) for 6 students:

\[
\begin{array}{c|cccccc}
x & 2 & 4 & 5 & 7 & 9 & 11 \\
\hline
y & 45 & 62 & 68 & 78 & 85 & 92
\end{array}
\]

\begin{enumerate}[label=(\alph*)]
\item Compute the OLS estimates \(\hat{\beta}_0\) and \(\hat{\beta}_1\) using formulas:
\[
\hat{\beta}_1 = \frac{\sum(x_i - \bar{x})(y_i - \bar{y})}{\sum(x_i - \bar{x})^2}, \quad \hat{\beta}_0 = \bar{y} - \hat{\beta}_1\bar{x}
\]

\item Calculate the residuals \(e_i = y_i - \hat{y}_i\), sum of squared errors (SSE), total sum of squares (SST), and coefficient of determination \(R^2\).

\item Compute the residual standard error \(s = \sqrt{\text{SSE}/(n-2)}\) and test the null hypothesis \(H_0: \beta_1 = 0\) at \(\alpha = 0.05\). Use the test statistic:
\[
t = \frac{\hat{\beta}_1 - 0}{s/\sqrt{\sum(x_i - \bar{x})^2}}
\]

\item Construct a 95\% confidence interval for \(\beta_1\).
\end{enumerate}

\subsection*{Problem 2: Multiple Linear Regression Matrix Formulation}

Consider a real estate dataset with house prices (\(y\), in thousands) predicted by area (\(x_1\), in 100 sq ft) and number of bedrooms (\(x_2\)):

\[
X = \begin{pmatrix}
1 & 12 & 2 \\
1 & 15 & 3 \\
1 & 18 & 3 \\
1 & 20 & 4 \\
1 & 25 & 4
\end{pmatrix}, \quad
\mathbf{y} = \begin{pmatrix} 180 \\ 220 \\ 260 \\ 280 \\ 340 \end{pmatrix}
\]

\begin{enumerate}[label=(\alph*)]
\item Compute \(X^TX\) and \(X^T\mathbf{y}\).

\item Calculate the OLS solution using the normal equations:
\[
\hat{\beta} = (X^TX)^{-1}X^T\mathbf{y}
\]

\item Compute the fitted values \(\hat{\mathbf{y}} = X\hat{\beta}\) and residuals \(\mathbf{e} = \mathbf{y} - \hat{\mathbf{y}}\).

\item Calculate \(R^2 = 1 - \frac{\text{SSE}}{\text{SST}}\) where \(\text{SSE} = \mathbf{e}^T\mathbf{e}\) and \(\text{SST} = \sum(y_i - \bar{y})^2\).
\end{enumerate}

\subsection*{Problem 3: Linear Regression with Categorical Predictors}

A company analyzes employee salaries (\(y\), in thousands) based on years of experience (\(x_1\)) and department (categorical: Sales, Tech, HR). Data for 8 employees:

\[
\begin{array}{c|c|c|c}
\text{Employee} & \text{Experience } (x_1) & \text{Department} & \text{Salary } (y) \\
\hline
1 & 2 & \text{Sales} & 45 \\
2 & 5 & \text{Sales} & 58 \\
3 & 3 & \text{Tech} & 62 \\
4 & 7 & \text{Tech} & 78 \\
5 & 4 & \text{Tech} & 68 \\
6 & 2 & \text{HR} & 42 \\
7 & 6 & \text{HR} & 55 \\
8 & 8 & \text{Sales} & 72
\end{array}
\]

\begin{enumerate}[label=(\alph*)]
\item Create dummy variables for the department (use Sales as reference category):
\[
D_{\text{Tech}} = \begin{cases} 1 & \text{if Tech} \\ 0 & \text{otherwise} \end{cases}, \quad
D_{\text{HR}} = \begin{cases} 1 & \text{if HR} \\ 0 & \text{otherwise} \end{cases}
\]
Write out the design matrix \(X\) (including intercept).

\item Set up and solve the normal equations to find \(\hat{\beta} = (\hat{\beta}_0, \hat{\beta}_1, \hat{\beta}_2, \hat{\beta}_3)^T\) where:
\[
y = \beta_0 + \beta_1 x_1 + \beta_2 D_{\text{Tech}} + \beta_3 D_{\text{HR}} + \epsilon
\]

\item Interpret each coefficient: What does \(\hat{\beta}_2\) represent?

\item Predict the salary for a Tech employee with 5 years of experience.
\end{enumerate}

\subsection*{Problem 4: Residual Analysis and Assumptions}

Given the simple linear model \(y = \beta_0 + \beta_1 x + \epsilon\) with fitted values and residuals:

\[
\begin{array}{c|c|c|c}
x & y & \hat{y} & e \\
\hline
1 & 3.2 & 2.8 & 0.4 \\
2 & 4.5 & 4.6 & -0.1 \\
3 & 6.8 & 6.4 & 0.4 \\
4 & 8.1 & 8.2 & -0.1 \\
5 & 10.5 & 10.0 & 0.5 \\
6 & 11.8 & 11.8 & 0.0 \\
7 & 14.2 & 13.6 & 0.6 \\
8 & 15.1 & 15.4 & -0.3
\end{array}
\]

\begin{enumerate}[label=(\alph*)]
\item Compute the mean and variance of residuals. Verify that \(\sum e_i = 0\).

\item Calculate standardized residuals: \(r_i = e_i / s\) where \(s = \sqrt{\sum e_i^2/(n-2)}\).

\item Check for outliers: Are any \(|r_i| > 2\)?

\item Compute the Durbin-Watson statistic to test for autocorrelation:
\[
DW = \frac{\sum_{i=2}^{n}(e_i - e_{i-1})^2}{\sum_{i=1}^{n}e_i^2}
\]
Values near 2 indicate no autocorrelation. Interpret your result.
\end{enumerate}

\subsection*{Problem 5: Weighted Least Squares}

When error variance is not constant (\(\text{Var}(\epsilon_i) = \sigma^2/w_i\)), we use weighted least squares (WLS). Consider data with known weights:

\[
\begin{array}{c|c|c}
x & y & w_i \\
\hline
1 & 2.5 & 1 \\
2 & 3.8 & 2 \\
3 & 5.5 & 2 \\
4 & 6.8 & 4 \\
5 & 8.2 & 4
\end{array}
\]

\begin{enumerate}[label=(\alph*)]
\item The WLS estimator minimizes \(\sum w_i(y_i - \beta_0 - \beta_1 x_i)^2\). Show that the WLS solution is:
\[
\hat{\beta}_1^{WLS} = \frac{\sum w_i(x_i - \bar{x}_w)(y_i - \bar{y}_w)}{\sum w_i(x_i - \bar{x}_w)^2}
\]
where \(\bar{x}_w = \sum w_i x_i / \sum w_i\).

\item Compute \(\bar{x}_w\), \(\bar{y}_w\), and the weighted sums needed for \(\hat{\beta}_1^{WLS}\) and \(\hat{\beta}_0^{WLS}\).

\item Compare the WLS estimates with OLS (unweighted) estimates. Which should be preferred?

\item Calculate the weighted residuals \(e_i^w = \sqrt{w_i}(y_i - \hat{y}_i)\) and verify the WLS fit quality.
\end{enumerate}

% ============================================================
% LOGISTIC REGRESSION PROBLEMS
% ============================================================

\section{Logistic Regression}

\subsection*{Problem 6: Binary Logistic Regression Fundamentals}

A credit card company models default probability based on debt-to-income ratio (\(x\)). Data for 8 customers:

\[
\begin{array}{c|c|c}
\text{Customer} & \text{Debt-to-Income } (x) & \text{Default } (y) \\
\hline
1 & 0.15 & 0 \\
2 & 0.25 & 0 \\
3 & 0.35 & 0 \\
4 & 0.40 & 1 \\
5 & 0.50 & 1 \\
6 & 0.60 & 1 \\
7 & 0.70 & 1 \\
8 & 0.80 & 1
\end{array}
\]

\begin{enumerate}[label=(\alph*)]
\item Write the logistic regression model:
\[
\log\left(\frac{p}{1-p}\right) = \beta_0 + \beta_1 x
\]
where \(p = P(y = 1 | x)\). Express \(p\) in terms of the logit function.

\item Suppose the estimated coefficients are \(\hat{\beta}_0 = -2.5\) and \(\hat{\beta}_1 = 8\). Calculate the predicted probability of default for \(x = 0.45\).

\item Find the decision boundary: At what debt-to-income ratio is \(p = 0.5\)?

\item Compute the odds ratio for a 0.1 unit increase in debt-to-income ratio:
\[
OR = \frac{p(x + 0.1) / (1 - p(x + 0.1))}{p(x) / (1 - p(x))} = e^{0.1\beta_1}
\]
Interpret this value.
\end{enumerate}

\subsection*{Problem 7: Maximum Likelihood Estimation for Logistic Regression}

For the logistic model with one predictor, the likelihood function is:
\[
L(\beta_0, \beta_1) = \prod_{i=1}^{n} p_i^{y_i}(1-p_i)^{1-y_i}
\]
where \(p_i = \frac{e^{\beta_0 + \beta_1 x_i}}{1 + e^{\beta_0 + \beta_1 x_i}}\).

Consider simplified data with 3 observations:

\[
\begin{array}{c|c|c}
x & y & \text{Group} \\
\hline
1 & 0 & A \\
2 & 1 & B \\
3 & 1 & C
\end{array}
\]

\begin{enumerate}[label=(\alph*)]
\item Write the log-likelihood function:
\[
\ell(\beta_0, \beta_1) = \sum_{i=1}^{n} \left[y_i(\beta_0 + \beta_1 x_i) - \log(1 + e^{\beta_0 + \beta_1 x_i})\right]
\]

\item Derive the gradient (score equations):
\[
\frac{\partial \ell}{\partial \beta_0} = \sum_{i=1}^{n}(y_i - p_i), \quad \frac{\partial \ell}{\partial \beta_1} = \sum_{i=1}^{n}x_i(y_i - p_i)
\]

\item For \(\beta_0 = 0\) and \(\beta_1 = 1\), evaluate the log-likelihood \(\ell(0, 1)\) and the gradient components.

\item Explain why we need iterative methods (Newton-Raphson or gradient descent) to find MLEs.
\end{enumerate}

\subsection*{Problem 8: Multinomial Logistic Regression}

A marketing study predicts customer preference (Low, Medium, High) based on price sensitivity (\(x\)). Use High as the reference category. The model is:

\[
\log\left(\frac{P(y = \text{Low})}{P(y = \text{High})}\right) = \beta_{01} + \beta_{11}x
\]
\[
\log\left(\frac{P(y = \text{Medium})}{P(y = \text{High})}\right) = \beta_{02} + \beta_{12}x
\]

Given: \(\beta_{01} = 2\), \(\beta_{11} = -3\), \(\beta_{02} = 1\), \(\beta_{12} = -1.5\).

\begin{enumerate}[label=(\alph*)]
\item For a customer with \(x = 0.5\), compute the linear predictors:
\[
\eta_1 = \beta_{01} + \beta_{11}(0.5), \quad \eta_2 = \beta_{02} + \beta_{12}(0.5)
\]

\item Calculate the probabilities using the softmax function:
\[
P(y = \text{Low}) = \frac{e^{\eta_1}}{1 + e^{\eta_1} + e^{\eta_2}}, \quad P(y = \text{Medium}) = \frac{e^{\eta_2}}{1 + e^{\eta_1} + e^{\eta_2}}, \quad P(y = \text{High}) = \frac{1}{1 + e^{\eta_1} + e^{\eta_2}}
\]

\item Verify that probabilities sum to 1.

\item Which category has the highest predicted probability for this customer?
\end{enumerate}

\subsection*{Problem 9: Model Evaluation with Confusion Matrix}

A logistic regression classifier predicts disease status (1 = diseased, 0 = healthy) for 10 patients. The threshold for classification is \(p \geq 0.5\):

\[
\begin{array}{c|c|c|c}
\text{Patient} & \text{True } y & \text{Predicted } p & \text{Predicted } \hat{y} \\
\hline
1 & 0 & 0.15 & 0 \\
2 & 0 & 0.30 & 0 \\
3 & 0 & 0.55 & 1 \\
4 & 1 & 0.45 & 0 \\
5 & 1 & 0.65 & 1 \\
6 & 1 & 0.75 & 1 \\
7 & 1 & 0.80 & 1 \\
8 & 0 & 0.20 & 0 \\
9 & 1 & 0.70 & 1 \\
10 & 0 & 0.40 & 0
\end{array}
\]

\begin{enumerate}[label=(\alph*)]
\item Construct the confusion matrix:
\[
\begin{array}{cc|cc}
 & & \multicolumn{2}{c}{\text{Predicted}} \\
 & & 0 & 1 \\
\hline
\text{Actual} & 0 & TN & FP \\
 & 1 & FN & TP
\end{array}
\]

\item Calculate accuracy, precision, recall (sensitivity), and specificity:
\[
\text{Accuracy} = \frac{TP + TN}{n}, \quad \text{Precision} = \frac{TP}{TP + FP}, \quad \text{Recall} = \frac{TP}{TP + FN}, \quad \text{Specificity} = \frac{TN}{TN + FP}
\]

\item Compute the F1-score:
\[
F1 = 2 \cdot \frac{\text{Precision} \times \text{Recall}}{\text{Precision} + \text{Recall}}
\]

\item If the threshold is changed to \(p \geq 0.6\), how does the confusion matrix change? Discuss the trade-off between sensitivity and specificity.
\end{enumerate}

\subsection*{Problem 10: ROC Curve and AUC}

For the data in Problem 9, we vary the classification threshold to generate the ROC curve.

\begin{enumerate}[label=(\alph*)]
\item For thresholds \(\tau = 0.4, 0.5, 0.6, 0.7\), compute the true positive rate (TPR) and false positive rate (FPR):
\[
TPR = \frac{TP}{TP + FN}, \quad FPR = \frac{FP}{FP + TN}
\]
Create a table of (FPR, TPR) pairs.

\item Plot the ROC curve conceptually (sketch the points and connect them).

\item The area under the ROC curve (AUC) can be approximated using the trapezoidal rule. For two consecutive points \((FPR_i, TPR_i)\) and \((FPR_{i+1}, TPR_{i+1})\):
\[
\text{Area}_i = \frac{1}{2}(FPR_{i+1} - FPR_i)(TPR_i + TPR_{i+1})
\]
Compute the total AUC.

\item Interpret the AUC value. What does AUC = 0.5 represent? What about AUC = 1.0?
\end{enumerate}

% ============================================================
% POLYNOMIAL REGRESSION PROBLEMS
% ============================================================

\section{Polynomial Regression}

\subsection*{Problem 11: Quadratic Regression Fitting}

A physicist models the height (\(y\), in meters) of a projectile as a function of time (\(x\), in seconds):

\[
\begin{array}{c|cccccc}
x & 0 & 0.5 & 1.0 & 1.5 & 2.0 & 2.5 \\
\hline
y & 0 & 11 & 18 & 21 & 20 & 15
\end{array}
\]

Model: \(y = \beta_0 + \beta_1 x + \beta_2 x^2 + \epsilon\)

\begin{enumerate}[label=(\alph*)]
\item Construct the design matrix \(X\):
\[
X = \begin{pmatrix}
1 & x_1 & x_1^2 \\
1 & x_2 & x_2^2 \\
\vdots & \vdots & \vdots \\
1 & x_6 & x_6^2
\end{pmatrix}
\]

\item Compute \(X^TX\) and \(X^T\mathbf{y}\).

\item Solve the normal equations to find \(\hat{\beta} = (X^TX)^{-1}X^T\mathbf{y}\).

\item Find the vertex of the parabola (maximum height): \(x^* = -\hat{\beta}_1/(2\hat{\beta}_2)\). Calculate the maximum height \(\hat{y}(x^*)\).

\item Compute \(R^2\) for the quadratic model. Compare conceptually with a linear model (without \(x^2\) term).
\end{enumerate}

\subsection*{Problem 12: Orthogonal Polynomials}

When fitting high-degree polynomials, multicollinearity among \(x, x^2, x^3, \ldots\) causes numerical instability. Orthogonal polynomials solve this.

Consider data: \(x = (-2, -1, 0, 1, 2)\).

\begin{enumerate}[label=(\alph*)]
\item Show that the raw polynomial design matrix columns are correlated. Compute the correlation between \(x\) and \(x^2\):
\[
\text{cor}(x, x^2) = \frac{\sum(x_i - \bar{x})(x_i^2 - \overline{x^2})}{\sqrt{\sum(x_i - \bar{x})^2 \sum(x_i^2 - \overline{x^2})^2}}
\]

\item Construct orthogonal polynomials using Gram-Schmidt:
\begin{itemize}
\item \(P_0(x) = 1\)
\item \(P_1(x) = x - \bar{x}\)
\item \(P_2(x) = x^2 - \langle x^2, P_0 \rangle P_0 - \langle x^2, P_1 \rangle P_1\)
\end{itemize}
where \(\langle u, v \rangle = \sum u_i v_i / n\).

\item Verify orthogonality: \(\sum P_1(x_i) P_2(x_i) = 0\).

\item Explain why using orthogonal polynomials improves numerical stability in regression.
\end{enumerate}

\subsection*{Problem 13: Overfitting in Polynomial Regression}

A data analyst fits polynomials of varying degrees to the data:

\[
\begin{array}{c|ccccc}
x & 1 & 2 & 3 & 4 & 5 \\
\hline
y & 2.1 & 3.9 & 6.2 & 7.8 & 10.1
\end{array}
\]

\begin{enumerate}[label=(\alph*)]
\item For a degree-1 polynomial (linear), the fitted equation is \(\hat{y} = 0.1 + 2.0x\) with \(R^2 = 0.985\).

For a degree-4 polynomial, \(R^2 = 1.000\) (perfect fit).

Compute the residual sum of squares (SSE) for both models. Which has smaller training error?

\item Calculate the adjusted \(R^2\) for both models:
\[
R^2_{\text{adj}} = 1 - \frac{(1-R^2)(n-1)}{n-p-1}
\]
where \(p\) is the number of predictors. Which model is preferred by adjusted \(R^2\)?

\item Explain the bias-variance trade-off: Why might the degree-4 polynomial perform poorly on new data despite perfect training fit?

\item Compute Akaike Information Criterion (AIC) for model selection:
\[
AIC = n \log(\text{SSE}/n) + 2p
\]
Use \(\text{SSE}_{\text{linear}} = 0.5\) and \(\text{SSE}_{\text{degree-4}} = 0.001\). Which model is preferred by AIC?
\end{enumerate}

\subsection*{Problem 14: Polynomial Regression with Interaction Terms}

A botanist studies plant growth (\(y\)) as a function of sunlight hours (\(x_1\)) and water amount (\(x_2\)). The model includes a quadratic term and interaction:

\[
y = \beta_0 + \beta_1 x_1 + \beta_2 x_2 + \beta_3 x_1^2 + \beta_4 x_1 x_2 + \epsilon
\]

Data for 5 plants:
\[
\begin{array}{c|c|c|c}
\text{Plant} & x_1 & x_2 & y \\
\hline
1 & 4 & 2 & 15 \\
2 & 6 & 3 & 25 \\
3 & 5 & 4 & 28 \\
4 & 8 & 2 & 30 \\
5 & 7 & 5 & 40
\end{array}
\]

\begin{enumerate}[label=(\alph*)]
\item Construct the design matrix \(X\) including all terms: intercept, \(x_1\), \(x_2\), \(x_1^2\), and \(x_1 x_2\).

\item Given the fitted coefficients: \(\hat{\beta} = (2, 3, 2.5, -0.2, 0.5)^T\), predict the growth for a plant with \(x_1 = 6\) and \(x_2 = 4\).

\item Interpret the interaction term \(\beta_4\): How does the effect of sunlight on growth depend on water amount?

\item Find the optimal sunlight hours (maximizing \(y\)) for \(x_2 = 3\) by taking \(\frac{\partial y}{\partial x_1} = 0\) and solving:
\[
\beta_1 + 2\beta_3 x_1 + \beta_4 x_2 = 0
\]
\end{enumerate}

\subsection*{Problem 15: Spline Regression as Piecewise Polynomials}

Spline regression fits different polynomials in different regions, connected smoothly at knots. Consider a cubic spline with one knot at \(x = 3\):

\[
f(x) = \begin{cases}
\beta_{01} + \beta_{11}x + \beta_{21}x^2 + \beta_{31}x^3 & \text{if } x \leq 3 \\
\beta_{02} + \beta_{12}x + \beta_{22}x^2 + \beta_{32}x^3 & \text{if } x > 3
\end{cases}
\]

For continuity and smoothness, we require:
\begin{itemize}
\item Continuity at \(x = 3\): \(f(3^-) = f(3^+)\)
\item First derivative continuity: \(f'(3^-) = f'(3^+)\)
\item Second derivative continuity: \(f''(3^-) = f''(3^+)\)
\end{itemize}

\begin{enumerate}[label=(\alph*)]
\item Write the three continuity conditions in terms of the coefficients.

\item Given \(\beta_{01} = 1\), \(\beta_{11} = 2\), \(\beta_{21} = -0.5\), \(\beta_{31} = 0.1\), and assuming the constraints are satisfied, compute \(f(3)\) and \(f'(3)\).

\item Explain why splines are preferred over high-degree polynomials for flexible curve fitting.

\item In practice, splines are often represented using basis functions. Write the truncated power basis function for a cubic spline:
\[
(x - c)_+^3 = \begin{cases} (x-c)^3 & \text{if } x > c \\ 0 & \text{otherwise} \end{cases}
\]
How does this represent a knot at \(x = c\)?
\end{enumerate}

% ============================================================
% RIDGE REGRESSION PROBLEMS
% ============================================================

\section{Ridge Regression}

\subsection*{Problem 16: Ridge Regression Derivation and Computation}

Consider the housing price model with severe multicollinearity:

\[
X = \begin{pmatrix}
1 & 1200 & 2.4 \\
1 & 1800 & 3.6 \\
1 & 2000 & 4.0 \\
1 & 2500 & 5.0
\end{pmatrix}, \quad
\mathbf{y} = \begin{pmatrix} 250 \\ 320 \\ 360 \\ 420 \end{pmatrix}
\]

Note: The second and third columns are perfectly correlated (\(x_2 = 0.002 x_1\)).

\begin{enumerate}[label=(\alph*)]
\item Compute \(X^TX\) and show that it has a very small eigenvalue (near-singular), indicating multicollinearity. Calculate the condition number \(\kappa = \lambda_{\max}/\lambda_{\min}\).

\item Derive the ridge regression estimator by minimizing:
\[
\|\mathbf{y} - X\beta\|^2 + \lambda\|\beta\|^2
\]
Show that the solution is:
\[
\hat{\beta}^{\text{ridge}} = (X^TX + \lambda I)^{-1}X^T\mathbf{y}
\]

\item For \(\lambda = 50\), compute \(X^TX + 50I\) and find \(\hat{\beta}^{\text{ridge}}\).

\item Compare with the OLS solution (if computable). Show that ridge shrinks coefficients toward zero.
\end{enumerate}

\subsection*{Problem 17: Bias-Variance Trade-off in Ridge Regression}

The mean squared error (MSE) of an estimator decomposes as:
\[
\text{MSE}(\hat{\beta}) = \text{Bias}^2(\hat{\beta}) + \text{Var}(\hat{\beta})
\]

For ridge regression, the bias and variance are:
\[
\text{Bias}(\hat{\beta}^{\text{ridge}}) = -\lambda(X^TX + \lambda I)^{-1}\beta
\]
\[
\text{Var}(\hat{\beta}^{\text{ridge}}) = \sigma^2 X(X^TX + \lambda I)^{-2}X^T
\]

\begin{enumerate}[label=(\alph*)]
\item Show that as \(\lambda \to 0\), ridge reduces to OLS: \(\hat{\beta}^{\text{ridge}} \to \hat{\beta}^{\text{OLS}}\). What happens to bias and variance?

\item Show that as \(\lambda \to \infty\), \(\hat{\beta}^{\text{ridge}} \to \mathbf{0}\). What happens to bias and variance?

\item For the simple case where \(X^TX = \text{diag}(\lambda_1, \lambda_2, \ldots, \lambda_p)\) (diagonal), show that the ridge shrinkage factor for the \(j\)-th coefficient is:
\[
\frac{\lambda_j}{\lambda_j + \lambda}
\]

\item Explain intuitively: Why does ridge regression improve prediction accuracy despite introducing bias?
\end{enumerate}

\subsection*{Problem 18: Cross-Validation for Lambda Selection}

To choose the optimal regularization parameter \(\lambda\), we use cross-validation. Consider Leave-One-Out CV (LOOCV):

Data: \(X = \begin{pmatrix} 1 & 2 \\ 1 & 4 \\ 1 & 6 \\ 1 & 8 \end{pmatrix}\), \(\mathbf{y} = \begin{pmatrix} 3 \\ 7 \\ 11 \\ 15 \end{pmatrix}\)

\begin{enumerate}[label=(\alph*)]
\item For \(\lambda = 1\), compute the ridge coefficients \(\hat{\beta}^{\text{ridge}}\) using all data.

\item For LOOCV, we leave out observation \(i\) and compute \(\hat{y}_{-i}\) (prediction for \(i\) using model trained on remaining data). The LOOCV shortcut formula is:
\[
\hat{y}_i - \hat{y}_{-i} = \frac{e_i}{1 - H_{ii}}
\]
where \(H = X(X^TX + \lambda I)^{-1}X^T\) is the hat matrix and \(e_i = y_i - \hat{y}_i\).

Compute the hat matrix \(H\) for \(\lambda = 1\).

\item Calculate the LOOCV error:
\[
CV(\lambda) = \frac{1}{n}\sum_{i=1}^{n}\left(\frac{e_i}{1 - H_{ii}}\right)^2
\]

\item Repeat for \(\lambda = 0.5\) and \(\lambda = 2\). Which \(\lambda\) minimizes CV error?
\end{enumerate}

\subsection*{Problem 19: Ridge Regression Path and Degrees of Freedom}

The ridge trace plots coefficient paths \(\hat{\beta}_j(\lambda)\) as \(\lambda\) varies. The effective degrees of freedom is:
\[
df(\lambda) = \text{tr}(H_\lambda) = \text{tr}(X(X^TX + \lambda I)^{-1}X^T)
\]

For the dataset:
\[
X = \begin{pmatrix} 1 & 1 \\ 1 & 2 \\ 1 & 3 \end{pmatrix}, \quad \mathbf{y} = \begin{pmatrix} 2 \\ 3 \\ 5 \end{pmatrix}
\]

\begin{enumerate}[label=(\alph*)]
\item Compute \(\hat{\beta}^{\text{ridge}}\) for \(\lambda = 0, 1, 5, 10\).

\item For each \(\lambda\), calculate \(df(\lambda) = \text{tr}(H_\lambda)\). Verify that \(df(0) = p\) (number of predictors) and \(df(\infty) \to 0\).

\item Plot the ridge path conceptually: Sketch how \(\hat{\beta}_1\) and \(\hat{\beta}_2\) change as \(\lambda\) increases from 0 to 10.

\item The generalized cross-validation (GCV) criterion is:
\[
GCV(\lambda) = \frac{\sum e_i^2}{(1 - df(\lambda)/n)^2}
\]
For \(\lambda = 1\), compute GCV if \(\sum e_i^2 = 0.5\) and \(df(1) = 1.5\).
\end{enumerate}

\subsection*{Problem 20: Ridge Regression vs LASSO}

Compare ridge regression with LASSO (Least Absolute Shrinkage and Selection Operator):

\textbf{Ridge:} \(\min_\beta \|\mathbf{y} - X\beta\|^2 + \lambda\|\beta\|_2^2\)

\textbf{LASSO:} \(\min_\beta \|\mathbf{y} - X\beta\|^2 + \lambda\|\beta\|_1\)

Consider a simple case with two predictors and OLS estimates \(\hat{\beta}^{\text{OLS}} = (6, 8)^T\).

\begin{enumerate}[label=(\alph*)]
\item Sketch the constraint regions in \((\beta_1, \beta_2)\) space:
\begin{itemize}
\item Ridge: \(\beta_1^2 + \beta_2^2 \leq t\) (circle)
\item LASSO: \(|\beta_1| + |\beta_2| \leq t\) (diamond)
\end{itemize}
Also sketch the OLS contours (ellipses centered at \(\hat{\beta}^{\text{OLS}}\)).

\item Explain geometrically why LASSO can produce exact zeros (sparse solutions) while ridge only shrinks toward zero.

\item For the constraint \(\beta_1^2 + \beta_2^2 = 25\), if ridge yields \(\hat{\beta}^{\text{ridge}} = (3, 4)^T\), find the corresponding \(\lambda\) by solving:
\[
\hat{\beta}^{\text{ridge}} = (X^TX + \lambda I)^{-1}X^T\mathbf{y}
\]
(Assume \(X^TX = I\) for simplicity and \(X^T\mathbf{y} = (6, 8)^T\)).

\item Discuss when to use ridge vs LASSO: Which is better for prediction? For interpretation?
\end{enumerate}

\end{document}
